\subsection{Reporte de limpieza de datos}

Hicimos la limpieza sobre los problemas que encontramos en 2.4. Para cada uno decidimos si lo corregimos, lo eliminamos o lo ignoramos, que son las tres estrategias del curso. Ningún archivo original se borra, ya que lo que se elimina solo deja de usarse y queda registrado en \nolinkurl{fotos.csv}, en las columnas \texttt{excluida} y \texttt{motivo}.

\begin{table}[H]
  \centering
  \small
  \caption{Problemas de calidad y cómo los tratamos.}
  \label{tab:limpieza}
  \begin{tabularx}{\textwidth}{>{\raggedright\arraybackslash}p{0.22\textwidth}>{\raggedright\arraybackslash}p{0.11\textwidth}>{\raggedright\arraybackslash}X>{\raggedright\arraybackslash}p{0.25\textwidth}}
    \toprule
    Problema (2.4) & Estrategia & Qué hicimos & Resultado \\
    \midrule
    Copias exactas dentro de una misma vaca (carpetas descargadas dos veces y nombres tipo \nolinkurl{IMG_0530(1).HEIC}) & Eliminar & Dejamos una sola copia por huella digital (SHA-256) dentro de cada vaca. Cuando hay varias, nos quedamos con la que no lleva ``(1)''. & De 761 archivos se quitaron 136 copias (122 de las carpetas del 14 de octubre y 14 de la vaca 91579) y quedaron 625 fotos únicas. \\
    \addlinespace
    Fotos idénticas en dos vacas (1548 y 2107) & Eliminar & Quitamos esas fotos de las dos vacas porque no sabemos a cuál pertenecen. & Se excluyeron 8 filas de fotos (4 fotos en 2 vacas). La 1548 conserva sus 3 fotos propias y la 2107 se queda sin fotos. \\
    \addlinespace
    Vaca 2123 con dos carpetas de fotos distintas que se enciman & Eliminar & Excluimos la vaca completa, porque no se sabe qué fotos van con cada una de sus dos evaluaciones y adivinar podría asignar un BCS equivocado. & Se excluyeron 14 fotos y 2 evaluaciones. \\
    \addlinespace
    Vacas con fotos pero sin BCS (1549 y 1572) & Eliminar & Quitamos sus fotos porque sin etiqueta no sirven para entrenar ni para validar. La regla aplica a cualquier vaca que no tenga BCS. & Se excluyeron 20 fotos. \\
    \addlinespace
    Fotos que muestran dos vacas (``1538 y 8744\_atras'') & Eliminar & Quitamos esas fotos porque no se sabe a cuál vaca corresponden. La 1538 conserva sus demás fotos. & Se excluyeron 3 fotos. \\
    \addlinespace
    Hora de la bitácora con ``x'' o con otro formato & Corregir e ignorar & ``13 00'' pasa a ``13:00''. Las ``x'' y los formatos raros quedan vacíos y el valor original se guarda en \texttt{hora\_original}. La hora no se usa en el modelo. & De 110 evaluaciones, 43 quedan con la hora vacía, que son 14 con ``x'', 2 con otro formato y las 27 de la hoja, que no registra la hora. \\
    \addlinespace
    BCS con coma decimal o fuera de escala & Corregir y verificar & El BCS se convierte a número y se revisa que esté entre 1 y 5 en pasos de 0.25. & Ningún BCS quedó sin convertir, fuera de rango ni fuera de los pasos de 0.25. \\
    \addlinespace
    Fotos JPG con orientación EXIF distinta de 1 (Samsung S22 e iPhone 17 Pro Max de la primera visita) & Corregir & Los JPG con orientación distinta de 1 se enderezan. En los HEIC la librería ya aplica la rotación al abrirlos. & 57 JPG del dataset final se enderezaron, 20 del Samsung y 37 del iPhone 17 Pro Max. Los demás JPG ya estaban derechos y se copiaron sin cambios. \\
    \addlinespace
    Dos formatos (HEIC y JPG) y varias resoluciones & Corregir & Todo pasa a JPG (véase 3.5). La resolución no se cambia en esta fase. & Véase la sección~\ref{sec:reformateados}. \\
    \addlinespace
    28 vacas con dos o más evaluaciones de BCS & Ignorar, documentado & No eliminamos ninguna evaluación. Las guardamos todas y entrenamos con el promedio (véase 3.3.1). & 27 de esas vacas siguen en el dataset porque la 2123 se excluyó, y 3 tienen tres evaluaciones (1543, 1567 y 1569). \\
    \addlinespace
    Valores faltantes & Ignorar, documentado & No imputamos nada, porque no tenemos campos numéricos con faltantes que estimar. & No aplica. \\
    \bottomrule
  \end{tabularx}
\end{table}

Con la limpieza, las cifras quedan de la siguiente forma.

\begin{itemize}
  \item Se leyeron 761 archivos de fotos, se quitaron 136 copias exactas y quedaron 625 fotos únicas. Se excluyeron 45 (14 de la vaca 2123, 8 filas repetidas entre la 1548 y la 2107, 20 de vacas sin BCS y 3 que muestran dos vacas), así que siguen 580 fotos de 77 vacas.
  \item Hay 110 evaluaciones, 83 de la bitácora y 27 de la hoja. Se excluyeron las 2 de la vaca 2123, así que quedan 108 evaluaciones.
\end{itemize}

\subsubsection*{Pendientes que siguen abiertos y su posible efecto}

\begin{itemize}
  \item Hay clases de BCS con muy pocas vacas (2.0 viene de una sola vaca y 2.25 de dos), así que las métricas pueden moverse mucho con un solo error.
  \item Las vistas de las 580 fotos las marcó una sola persona del equipo y queda pendiente decidir si otra persona revisa una muestra de las marcas.
\end{itemize}
